\author{Brett Reynolds}
\title{Cảnh quan ngôn ngữ: Language Landscapes}
\subtitle{Hướng dẫn về cấu trúc tiếng Anh cho giáo viên tiếng Anh [Bản tiếng Việt]}
% \ISBNdigital{}
% \ISBNhardcover{}
% \BookDOI{}
\typesetter{Brett Reynolds}
% \proofreader{}
% \lsCoverTitleSizes{51.5pt}{17pt}
\renewcommand{\lsSeries}{tbls}
%\newcommand{\lsSeriesTitle}{Textbooks in Language Sciences}
%\newcommand{\lsSeriesColor}{lsYellow}
\renewcommand{\lsSeriesNumber}{16}
\BackBody{Ngôn ngữ không phải là một danh sách để học thuộc; đó là một cảnh quan để khảo sát.

Một nhà tự nhiên học giỏi ngoài thực địa không chỉ đặt tên cho sự vật. Họ quan sát chúng vận hành như thế nào, tụ lại ở đâu, và vì sao biến đổi. Cuốn sách này mời bạn đọc tiếng Anh theo cùng cách ấy: chú ý đến những gì người nói thật sự làm, chứ không chỉ những gì sách quy tắc bảo rằng họ nên làm. Khung phân tích dựa trên \textit{The Cambridge Grammar of the English Language}, được trình bày giản lược cho dễ tiếp cận nhưng không hạ thấp mức độ chính xác.

Các chương đi từ âm thanh và từ ngữ đến cụm từ, mệnh đề, rồi diễn ngôn, với những ví dụ tiếng Anh rõ ràng đặt cạnh các lát cắt nhỏ từ những ngôn ngữ khác. Mỗi chương luyện cùng một thao tác: quan sát một khuôn mẫu, kiểm tra nó bằng bằng chứng, nhận ra nơi khuôn mẫu ấy vỡ ra, rồi hỏi sự vỡ ra ấy dạy ta điều gì. Đến cuối sách, bạn sẽ có một vốn ngữ pháp thực dụng đủ sức bảo vệ các nhận định của mình: dùng để đánh giá sách giáo khoa, giải thích vì sao câu của người học nghe chưa ổn, và bảo vệ lời giải thích khi bị chất vấn.

Dù bạn đang trong chương trình đào tạo giáo viên hay đã đứng lớp, lợi ích vẫn như nhau: ngừng dựa vào những quy tắc không thể biện minh, và bắt đầu tự mình nhìn thấy hệ thống của ngôn ngữ.}
